\subsection{Familias de modelos}

La comparación incluye 28 candidatos individuales procedentes de tres familias y cinco reglas de ensemble. Todos producen cinco scores principales; cuando la arquitectura lo permite, añade 14 etiquetas finas y tres flags, para un total de 22 logits con pérdidas enmascaradas. La implementación usa scikit-learn para los modelos clásicos \cite{pedregosa2011sklearn} y PyTorch/Transformers para las redes \cite{paszke2019pytorch,wolf2020transformers}. La Fig.~\ref{fig:modelos} contrasta los mecanismos de las tres familias y anticipa su composición en el ensemble.

\begin{figure*}[t]
\centering
\resizebox{.98\textwidth}{!}{\begin{tikzpicture}[
  x=1cm,y=1cm,font=\scriptsize,
  box/.style={draw,rounded corners=2pt,text width=21mm,minimum height=7.5mm,align=center,inner sep=1.7pt},
  family/.style={box,fill=blue!7},
  map/.style={box,fill=purple!8,text width=23.5mm,minimum height=19mm},
  struct/.style={box,fill=orange!9},
  outputbox/.style={box,fill=green!10,text width=23.5mm,minimum height=13.5mm},
  arr/.style={-{Stealth[length=1.8mm]},thick},
  bus/.style={thick}
]
\node[box] (text) at (0,0) {Chunk textual};
\node[family] (classic) at (3.2,2.1) {TF--IDF\\SVM y logística};
\node[family] (minilm) at (3.2,.7) {MiniLM\\encoder de tokens};
\node[family] (e5) at (3.2,-.7) {E5\\encoder contrastivo};
\node[family] (qwen) at (3.2,-2.1) {Qwen3--0.6B\\LLM con LoRA};

\coordinate (inbus) at (1.55,0);
\draw[bus] (text.east)--(inbus);
\draw[bus] (inbus |- classic.west)--(inbus |- qwen.west);
\foreach \m in {classic,minilm,e5,qwen}{
  \draw[arr] (inbus |- \m.west)--(\m.west);
}

\node[map] (tested) at (6.8,0) {\textbf{Diseños probados}\\Plano\\Cascada\\Jerarquía o multitarea};
\coordinate (familybus) at (5.1,0);
\foreach \m in {classic,minilm,e5,qwen}{
  \draw[bus] (\m.east)--(familybus |- \m.east);
}
\draw[bus] (familybus |- classic.east)--(familybus |- qwen.east);
\draw[arr] (familybus |- tested.west)--(tested.west);

\node[struct] (flat) at (10.5,1.5) {Plano\\4 categorías};
\node[struct] (cascade) at (10.5,0) {Cascada\\daño $\rightarrow$ categorías};
\node[struct] (hier) at (10.5,-1.5) {Jerarquía compartida\\o multitarea};
\coordinate (structbus) at (8.65,0);
\draw[bus] (tested.east)--(structbus);
\draw[bus] (structbus |- flat.west)--(structbus |- hier.west);
\foreach \s in {flat,cascade,hier}{
  \draw[arr] (structbus |- \s.west)--(\s.west);
}

\node[outputbox] (scores) at (14.2,0) {Scores\\umbrales por clase\\calibración indicada};
\coordinate (scorebus) at (12.25,0);
\foreach \s in {flat,cascade,hier}{
  \draw[bus] (\s.east)--(scorebus |- \s.east);
}
\draw[bus] (scorebus |- flat.east)--(scorebus |- hier.east);
\draw[arr] (scorebus |- scores.west)--(scores.west);
\end{tikzpicture}
}
\caption{Familias comparadas y composición del ensemble congelado. Los tres miembros conservan arquitecturas y errores diferentes. Elaboración propia.}
\label{fig:modelos}
\end{figure*}

\subsubsection{Modelos clásicos}

La vista base concatena TF--IDF de palabras (1--2 gramas) y caracteres internos (3--5 gramas), siguiendo la ponderación de frecuencia inversa de documento \cite{sparckjones1972idf}. Una transformación uno-contra-resto enmascarada entrena cabezas con \emph{Dummy}, Complement Naive Bayes, regresión logística, SVM lineal calibrada o SGD logístico \cite{cox1958logistic,cortes1995svm,rennie2003cnb}. Una vista adicional incorpora indicadores léxicos auditables derivados de la política de etiquetado; continúa siendo aprendizaje supervisado y no \emph{prompting}.

El mejor clásico según el ranking común es una regresión logística con $C=0,5$. Su interés no se limita al costo: aporta una señal léxica distinta a las redes y permite inspeccionar n-gramas, por lo que resulta un miembro complementario del ensemble.

\subsubsection{Transformers compactos y cascada ganadora}

Los modelos planos usan \path{sentence-transformers/paraphrase-multilingual-MiniLM-L12-v2} o \path{intfloat/multilingual-e5-small} \cite{wang2020minilm,wang2022e5,wang2024e5,hf2026minilmcard,hf2026e5card}. Se ajustan con AdamW, entropía cruzada binaria ponderada y parada temprana \cite{loshchilov2019adamw,pytorch2026bce,prechelt1998earlystopping}. También se evaluaron una cascada suave, una cabeza denominada multitarea y una cascada v2 orientada a seguridad.

La mejor configuración Transformer es la cascada v2. Una primera red E5 aprende \texttt{ANY\_DAMAGE} y auxiliares; una segunda red E5 produce \texttt{SEGURO} y los cuatro daños. Sea $g(x)$ el score de daño de la compuerta y $\tau$ su umbral. La compuerta declara seguro solo cuando $g(x)<\tau$; en otro caso usa los cinco scores de la rama. En validation se escoge el mayor $\tau$ que cumple simultáneamente

\[
\begin{aligned}
\operatorname{Recall}_{\text{daño}}(\tau)&\ge 0,99,\\
P(y=\texttt{SEGURO}\mid g(x)<\tau)&\ge 0,99.
\end{aligned}
\]

Si ningún umbral con cobertura positiva satisface ambas condiciones, todo pasa a la rama. La arquitectura separa detección amplia y tipificación, una decisión relacionada con antecedentes de clasificación jerárquica \cite{silla2011hierarchical}. Sin embargo, ejecuta dos Transformers y puede propagar errores de la compuerta; sus valores 0,99 describen validation y no una garantía poblacional.

\subsubsection{Qwen3--LoRA}

El estudio usa \path{Qwen/Qwen3-0.6B-Base}, un modelo denso multilingüe de 0,6 mil millones de parámetros \cite{qwen2025qwen3,hf2026qwen06bcard}. LoRA congela el backbone e inserta matrices entrenables de bajo rango \cite{hu2022lora}. La configuración ganadora coloca adaptadores en \texttt{q\_proj}, \texttt{k\_proj}, \texttt{v\_proj} y \texttt{o\_proj}, con rango 8, $\alpha=16$, \emph{dropout} 0,05 y tasa $10^{-4}$, hasta cuatro épocas. Una cabeza de 22 logits produce las cinco salidas principales y auxiliares. Este Qwen es un clasificador: no recibe la taxonomía Markdown ni genera JSON durante la inferencia.

Qwen3--LoRA es el mejor candidato individual de la comparación. Su representación más amplia aporta contexto; la adaptación eficiente evita ajustar todos los parámetros. El costo de memoria y latencia es superior al de regresión logística y E5, razón adicional para conservar las tres familias en el artefacto.

\subsection{Calibración y predicciones fuera de pliegue}

Cada candidato produce scores sobre validation mediante cinco pliegues \texttt{GroupKFold} por video \cite{sklearn2026groupkfold}. En cada pliegue, cuatro quintos ajustan calibradores sigmoidales de Platt y umbrales; el quinto produce predicciones fuera de pliegue (OOF) \cite{platt1999probabilistic,niculescu2005probabilities}. Así, la exactitud balanceada y la macro-AUPRC OOF no evalúan una fila con un calibrador ajustado sobre esa misma fila. Después de seleccionar, los calibradores y umbrales de despliegue se ajustan con toda validation y quedan congelados.

Para una categoría $k$, la precisión promedio resume la curva precisión--recobrado,
\(
\AP_k=\sum_n(R_n-R_{n-1})P_n
\), y el promedio macro asigna igual peso a los cuatro daños \cite{sklearn2026averageprecision,saito2015pr}. La compuerta \texttt{ANY\_DAMAGE} se evalúa con exactitud balanceada (BA), media de sensibilidad y especificidad, para no favorecer la clase segura mayoritaria.

\subsection{Elección e interpretación de las métricas}

Las métricas responden a decisiones distintas. La métrica principal del ranking es la BA de \texttt{ANY\_DAMAGE} OOF a cobertura completa: mide por igual la capacidad de detectar daño y reconocer contenido seguro. Una BA de 0,84 no significa que 84\% de todos los chunks sean correctos, sino que el promedio entre sensibilidad y especificidad es 0,84 \cite{brodersen2010balanced}. La macro-AUPRC promedia sin ponderación la precisión promedio de los cuatro daños a través de los umbrales; un valor de 0,55 no es un porcentaje de aciertos y solo se compara directamente sobre el mismo panel y prevalencia. La macro-F1 promedia por daño $2PR/(P+R)$ después de fijar umbrales y describe el compromiso entre precisión y recobrado, sin usar verdaderos negativos \cite{sokolova2009metrics}.

La exactitud ordinaria y los promedios micro o ponderados no gobiernan el ranking porque la abundancia de \texttt{SEGURO} y de las etiquetas con mayor soporte puede ocultar omisiones en daños raros. ROC-AUC se conserva como diagnóstico secundario; ante positivos escasos, la curva precisión--recobrado muestra de forma más directa cuántas alertas relevantes aparecen entre las predicciones positivas \cite{saito2015pr,davis2006pr}. Tampoco se suman BA, FNR, FPR, AUPRC y F1 en una nota única: varias cantidades comparten errores y sus escalas no representan la misma utilidad. El riesgo $R_{0,67}=0,67\,\mathrm{FNR}+0,33\,\mathrm{FPR}$ expresa la decisión local de penalizar aproximadamente el doble una omisión que una falsa alerta, mientras la tasa de revisión mide capacidad humana y se interpreta junto con el error de la ruta automática.

\subsection{Ensemble por promedio de probabilidades}

El ensemble congelado combina los mejores representantes de cada familia: regresión logística $C=0,5$, cascada E5 v2 y Qwen3--LoRA. Para cada salida $k$, promedia los tres scores sin pesos y calibra el promedio:

\[
\bar{s}_k(x)=\frac{1}{3}\sum_{m=1}^{3}s_{mk}(x),
\qquad
\hat p_k(x)=\sigma(a_k\bar{s}_k(x)+b_k).
\]

La decisión compara $\hat p_k$ con un umbral $\theta_k$ fijado en validation. Esta fusión suave aprovecha que las señales léxicas, contextuales y autoregresivas no fallan siempre en los mismos casos; los ensembles y el voto mayoritario cuentan con antecedentes generales \cite{dietterich2000ensemble,lam1997majority}. El experimento también evalúa promedio ponderado por validation, mayoría dura, unión (máximo) e intersección (mínimo). Solo el promedio simple se congela para test. El Apéndice~\ref{app:ensemble} conserva los identificadores, calibradores, umbrales y demás parámetros necesarios para reconstruirlo.

\subsection{Criterio de selección e inferencia estadística}

La política de selección es lexicográfica y no suma métricas redundantes. Primero maximiza BA \texttt{ANY\_DAMAGE} OOF a cobertura completa. Luego exige una salvaguarda de macro-AUPRC de daños frente al mejor individuo, con margen de no inferioridad predeclarado de 0,10; conserva la frontera de Pareto BA--macro-AUPRC; y desempata por menor riesgo $R_{0,67}=0,67\,\mathrm{FNR}+0,33\,\mathrm{FPR}$, seguido de mayor macro-AUPRC. El margen 0,10 es permisivo y se conserva como limitación, no como prueba de equivalencia.

La aplicación al resultado final muestra la función de cada paso. El promedio simple obtiene BA 0,840, frente a 0,838 del promedio ponderado y 0,831 de Qwen3--LoRA; por ello lidera el criterio primario. Su macro-AUPRC es 0,555, superior en 0,039 a 0,516 del mejor individuo, de modo que supera la salvaguarda. El promedio ponderado alcanza AUPRC 0,556 pero BA menor, por lo que ambos ensembles permanecen en la frontera de Pareto: ninguno mejora simultáneamente las dos coordenadas. Si se requiere el desempate de costo, el promedio simple tiene $R_{0,67}=0,67(0,127)+0,33(0,193)=0,149$, menor que 0,152 del ponderado. Estas operaciones se aplican en secuencia; no se suman en una puntuación compuesta.

Las diferencias se estiman con 2\,000 réplicas bootstrap pareadas por video y ajuste de Holm \cite{efron1979bootstrap,field2007clusterbootstrap,holm1979multiple}. El remuestreo respeta la correlación entre chunks de un mismo video; no sustituye la variación entre semillas.

\subsection{Política selectiva y test sellado}

\texttt{NEEDS\_REVIEW} es una política posterior al modelo. Se activa por conflicto seguro--daño, salida vacía, incoherencia compuerta--categoría o proximidad a los umbrales. La capacidad máxima declarada fue 40\% y validation eligió un margen $\delta=0,03$, con 33\% de revisión. Candidato, miembros, calibradores, umbrales y política se almacenaron en \texttt{seleccion\_congelada.json} antes de test.

El test natural contiene 22\,684 chunks y se abrió una sola vez. Los tres miembros se infirieron en esa apertura y sus scores se reutilizaron para calcular la vista secundaria 4:1; no hubo reinferencia, reajuste ni cambio de ganador. Este diseño separa selección y estimación final, y reduce el sesgo de selección discutido por Cawley y Talbot \cite{cawley2010selection}.

\subsection{Artefacto de demostración}

El frente de producción carga el registro de los tres ganadores por tipo y del ensemble. Acepta texto o URL con subtítulos, genera chunks temporales, permite elegir el modelo y muestra scores, categorías, razones de revisión y procedencia. Las acciones humanas se registran para auditoría y para una futura versión de entrenamiento separada. La Fig.~\ref{fig:deploy} resume esta arquitectura y el Apéndice~\ref{app:interfaces} documenta las interfaces. El demostrador opera localmente o en modo sombra; no está autorizado para afectar usuarios sin revisión.

\begin{figure*}[t]
\centering
\resizebox{\textwidth}{!}{\begin{tikzpicture}[
  font=\footnotesize,
  node distance=5mm and 6mm,
  box/.style={draw,rounded corners=2pt,text width=25mm,minimum height=10mm,align=center},
  input/.style={box,fill=blue!7},
  model/.style={box,fill=orange!9},
  human/.style={box,fill=green!10},
  store/.style={box,fill=purple!8},
  arr/.style={-{Stealth[length=2mm]},thick}
]
\node[input] (entry) {Una caja\\texto o URL};
\node[input,right=of entry] (detect) {Detección automática\\frase / YouTube};
\node[input,right=of detect] (prep) {Subtítulos exigidos\\chunks + tiempos};
\node[model,right=of prep,yshift=14mm] (svm) {SVM plano};
\node[model,right=of prep] (e5) {E5 plano};
\node[model,right=of prep,yshift=-14mm] (qwen) {Qwen3--LoRA\\plano};
\node[model,right=10mm of e5] (rule) {Modelo único,\\comparación o\\consenso 2 de 3};
\node[human,right=of rule] (decision) {Alerta + evidencia\\aceptar, rechazar\\o modificar};
\node[store,below=8mm of decision] (store) {Base y registro versionado\\modelo, categoría,\\decisión y versión};
\draw[arr] (entry)--(detect);
\draw[arr] (detect)--(prep);
% Dos buses en carriles vacíos sustituyen las diagonales de entrada y salida.
\coordinate (inbus) at ([xshift=3mm]prep.east);
\draw[thick] (prep.east)--(inbus);
\draw[thick] (inbus |- svm.west)--(inbus |- qwen.west);
\foreach \m in {svm,e5,qwen}{
  \draw[arr] (inbus |- \m.west)--(\m.west);
}
\coordinate (outbus) at ([xshift=-3mm]rule.west);
\foreach \m in {svm,e5,qwen}{
  \draw[thick] (\m.east)--(outbus |- \m.east);
}
\draw[thick] (outbus |- svm.east)--(outbus |- qwen.east);
\draw[arr] (outbus |- rule.west)--(rule.west);
\draw[arr] (rule)--(decision);
\draw[arr] (decision)--(store);
\draw[arr,dashed] (store.south) -- ++(0,-5mm) -| node[pos=.5,below]{futura iteración} (prep.south);
\end{tikzpicture}
}
\caption{Arquitectura operacional vigente: tres modelos por tipo, ensemble suave, política selectiva y decisión humana. Elaboración propia.}
\label{fig:deploy}
\end{figure*}
